\documentclass{article}
\usepackage{amsmath,amssymb,amsthm}
\usepackage{algorithm}
\usepackage{algorithmic}
\usepackage{bm}
\usepackage{hyperref}
\usepackage{enumitem}
\newtheorem{theorem}{Theorem}
\newtheorem{lemma}{Lemma}
\newtheorem{assumption}{Assumption}
\newtheorem{corollary}{Corollary}
\newcommand{\E}{\mathbb{E}}
\newcommand{\R}{\mathbb{R}}

\begin{document}

\section*{Adaptive Variance‑Weighted Gradient Descent (AVGD)}

\section*{Mathematical Formulation}
Let $f:\R^{d}\to\R$ be the objective function.  
Denote by $\nabla f(x)$ the (possibly stochastic) gradient evaluated at $x\in\R^{d}$.  
AVGD maintains an exponential moving average (EMA) of the gradients,
\begin{equation}
    v_{t}= (1-\alpha)v_{t-1}+ \alpha \nabla f(x_{t}),\qquad v_{0}=0,
    \label{eq:ema}
\end{equation}
with smoothing parameter $\alpha\in(0,1)$.  
Because $v_{t}$ is biased towards zero at early iterations, a bias–corrected estimator is used:
\begin{equation}
    \widehat v_{t}= \frac{v_{t}}{1-(1-\alpha)^{t}} .
    \label{eq:bias}
\end{equation}
The step size at iteration $t$ is adapted to the magnitude of $\widehat v_{t}$,
\begin{equation}
    \eta_{t}= \frac{\eta}{1+\gamma\|\widehat v_{t}\|},
    \label{eq:stepsize}
\end{equation}
where $\eta>0$ is a base learning rate and $\gamma\ge 0$ controls the sensitivity to the variance estimator.  
The parameter update is
\begin{equation}
    x_{t+1}= x_{t} - \eta_{t}\,\nabla f(x_{t}).
    \label{eq:update}
\end{equation}
Algorithm terminates when $\|\nabla f(x_{t})\|\le\epsilon$ or after a maximal number $T$ of iterations.

\section*{Pseudocode}
\begin{algorithm}[H]
\caption{Adaptive Variance‑Weighted Gradient Descent (AVGD)}
\begin{algorithmic}[1]
\REQUIRE Initial point $x_{0}\in\R^{d}$, base step size $\eta>0$, EMA parameter $\alpha\in(0,1)$,
        sensitivity $\gamma\ge0$, tolerance $\epsilon>0$, maximal iter. $T$.
\STATE $v_{0}\gets 0$, $t\gets 0$.
\WHILE{$t<T$ \textbf{and} $\|\nabla f(x_{t})\|>\epsilon$}
    \STATE Compute (stochastic) gradient $g_{t}\gets \nabla f(x_{t})$.
    \STATE Update EMA: $v_{t}= (1-\alpha)v_{t-1}+ \alpha g_{t}$.
    \STATE Bias correction: $\widehat v_{t}= v_{t}/\bigl(1-(1-\alpha)^{t}\bigr)$.
    \STATE Adaptive step size: $\displaystyle \eta_{t}= \frac{\eta}{1+\gamma\|\widehat v_{t}\|}$.
    \STATE Parameter update: $x_{t+1}= x_{t}-\eta_{t} g_{t}$.
    \STATE $t\gets t+1$.
\ENDWHILE
\RETURN $x_{t}$.
\end{algorithmic}
\end{algorithm}

\section*{Dry Run Example}
Consider $f(w)=(w-3)^{2}$, thus $\nabla f(w)=2(w-3)$.  
Choose $w_{0}=0$, $\eta=0.1$, $\alpha=0.5$, $\gamma=0$ (so $\eta_{t}=0.1$).  
The EMA recursion reduces to $v_{t}=0.5\,v_{t-1}+0.5\,\nabla f(w_{t})$ and bias correction is unnecessary because $\gamma=0$.

\begin{center}
\begin{tabular}{c|c|c|c|c}
$t$ & $w_{t}$ & $\nabla f(w_{t})$ & $v_{t}$ & $f(w_{t})$\\ \hline
0 & $0$ & $-6$ & $0$ & $9$\\ \hline
1 & $w_{1}=0-0.1(-6)=0.6$ &
   $2(0.6-3)=-4.8$ &
   $v_{1}=0.5\cdot0+0.5(-4.8)=-2.4$ &
   $f(0.6)=5.76$\\ \hline
2 & $w_{2}=0.6-0.1(-4.8)=1.08$ &
   $2(1.08-3)=-3.84$ &
   $v_{2}=0.5(-2.4)+0.5(-3.84)=-3.12$ &
   $f(1.08)=3.6864$\\ \hline
3 & $w_{3}=1.08-0.1(-3.84)=1.464$ &
   $2(1.464-3)=-3.072$ &
   $v_{3}=0.5(-3.12)+0.5(-3.072)=-3.096$ &
   $f(1.464)=2.3515$\\
\end{tabular}
\end{center}
The objective value decreases rapidly, illustrating the effect of the EMA‑based step‑size adaptation.

\newpage
\section*{Assumptions}
\begin{assumption}[Smoothness]\label{ass:smooth}
$f$ is $L$‑smooth, i.e.
\[
    \|\nabla f(x)-\nabla f(y)\|\le L\|x-y\|,\qquad\forall x,y\in\R^{d}.
\]
\end{assumption}
\begin{assumption}[Polyak–Łojasiewicz (PL) condition]\label{ass:pl}
There exists $\mu>0$ such that
\[
    \frac12\|\nabla f(x)\|^{2}\ge\mu\bigl(f(x)-f^{\star}\bigr),\qquad\forall x\in\R^{d},
\]
where $f^{\star}= \min_{x}f(x)$.
\end{assumption}
\begin{assumption}[Unbiased stochastic gradients]\label{ass:unbiased}
For each $t$, the gradient estimator $g_{t}=\nabla f(x_{t};\xi_{t})$ satisfies
\[
    \E\!\left[g_{t}\mid x_{t}\right]=\nabla f(x_{t}),\qquad
    \E\!\left[\|g_{t}-\nabla f(x_{t})\|^{2}\mid x_{t}\right]\le\sigma^{2},
\]
with a universal variance bound $\sigma^{2}\ge0$.
\end{assumption}
\begin{assumption}[Hyper‑parameter range]\label{ass:param}
\[
    0<\alpha<1,\qquad \eta\le\frac{1}{L},\qquad \gamma\ge0.
\]
\end{assumption}

\section*{Main Convergence Theorem}
\begin{theorem}\label{thm:linear}
Assume \ref{ass:smooth}–\ref{ass:param} hold. Let $\{x_{t}\}$ be generated by AVGD with the adaptive step size \eqref{eq:stepsize}. Define the contraction factor
\[
    \rho \;:=\; \frac{\mu\eta}{1+\gamma G},\qquad 
    G\;:=\;\sup_{t\ge0}\E\!\bigl[\|\widehat v_{t}\|\bigr].
\]
If $\rho\in(0,1)$, then for all $T\ge1$,
\[
    \E\!\bigl[f(x_{T})-f^{\star}\bigr]\;\le\;(1-\rho)^{T}\,\bigl(f(x_{0})-f^{\star}\bigr)
    \;+\; \frac{\eta\sigma^{2}}{2\mu(1+\gamma G)} .
\]
In particular, when $\sigma=0$ (exact gradients) the method enjoys a \emph{linear} convergence rate:
\[
    \E\!\bigl[f(x_{T})-f^{\star}\bigr]\le (1-\rho)^{T}\bigl(f(x_{0})-f^{\star}\bigr).
\]
\end{theorem}

\section*{Theoretical Properties}
\begin{itemize}[leftmargin=*]
    \item \textbf{Stability.} The EMA recursion \eqref{eq:ema} guarantees that $v_{t}$ remains bounded whenever the gradients are bounded; bias correction \eqref{eq:bias} removes the transient decay factor $(1-\alpha)^{t}$.
    \item \textbf{Hyper‑parameter sensitivity.} Larger $\alpha$ yields a variance estimator that tracks the current gradient more closely (smaller lag) but introduces higher variance; $\gamma$ controls how aggressively the step size shrinks when $\|\widehat v_{t}\|$ grows. The contraction factor $\rho$ decreases monotonically with $\gamma$.
    \item \textbf{Comparison to baselines.} Setting $\gamma=0$ recovers standard (constant‑step) gradient descent with rate $1-\mu\eta$, while $\gamma>0$ can improve robustness on noisy gradients because the effective step size is reduced when the EMA magnitude is large, similarly to Adam’s denominator but without element‑wise scaling.
\end{itemize}

\section*{Discussion}
The theorem shows that under the PL condition AVGD inherits the linear convergence of deterministic gradient descent, while the additional term $\frac{\eta\sigma^{2}}{2\mu(1+\gamma G)}$ quantifies the steady‑state error caused by stochastic noise. The bound on $G$ can be obtained from Lemma~\ref{lem:boundG} below.

\newpage
\section*{Preliminaries (Definitions and Lemmas)}
\begin{lemma}[Bound on the EMA norm]\label{lem:boundG}
Under Assumptions \ref{ass:smooth}–\ref{ass:unbiased} and $\alpha\in(0,1)$,
\[
    \E\!\bigl[\|v_{t}\|\bigr]\;\le\; \frac{\alpha}{1-\alpha}\,\bigl(L\,\E\!\|x_{t}-x^{\star}\|+\sigma\bigr),
\]
where $x^{\star}$ is any global minimizer. Consequently,
\[
    G\;\le\;\frac{\alpha}{1-\alpha}\,\bigl(L D+\sigma\bigr),\qquad 
    D:=\sup_{t}\E\!\bigl[\|x_{t}-x^{\star}\|\bigr] <\infty .
\]
\end{lemma}
\begin{proof}
[Step 1] Write the EMA recursion \eqref{eq:ema} and take norms:
\[
    \|v_{t}\|\le (1-\alpha)\|v_{t-1}\|+\alpha\|g_{t}\|.
\]
[Step 2] Apply expectation conditioned on the past and use the unbiasedness \eqref{ass:unbiased}:
\[
    \E\!\bigl[\|g_{t}\|\mid\mathcal{F}_{t-1}\bigr]
    \le \E\!\bigl[\|g_{t}-\nabla f(x_{t})\|\mid\mathcal{F}_{t-1}\bigr]
        +\|\nabla f(x_{t})\|
    \le \sigma+\|\nabla f(x_{t})\|.
\]
[Step 3] By $L$‑smoothness, $\|\nabla f(x_{t})\|\le L\|x_{t}-x^{\star}\|$.
[Step 4] Take total expectation and unroll the recursion:
\[
    \E\!\bigl[\|v_{t}\|\bigr]\le (1-\alpha)^{t}\|v_{0}\|
        +\alpha\sum_{k=0}^{t-1}(1-\alpha)^{k}\bigl(L\,\E\!\|x_{t-k}-x^{\star}\|+\sigma\bigr).
\]
[Step 5] Since $\|v_{0}\|=0$ and $(1-\alpha)^{k}\le1$, we obtain the claimed bound.
\end{proof}

\section*{Main Proof (Step‑by‑Step Derivation)}
We prove Theorem~\ref{thm:linear}. Throughout the proof, let $\mathcal{F}_{t}$ denote the $\sigma$‑algebra generated by $\{x_{0},\dots,x_{t}\}$.

\noindent\textbf{Step 1 (Smoothness inequality).}  
By $L$‑smoothness,
\begin{equation}
    f(x_{t+1})\le f(x_{t})+\langle\nabla f(x_{t}),x_{t+1}-x_{t}\rangle
        +\frac{L}{2}\|x_{t+1}-x_{t}\|^{2}.
    \label{eq:smooth}
\end{equation}
[Step 2] Insert the update \eqref{eq:update}:
\[
    x_{t+1}-x_{t}= -\eta_{t}g_{t}.
\]
Thus,
\[
    f(x_{t+1})\le f(x_{t}) -\eta_{t}\langle\nabla f(x_{t}),g_{t}\rangle
        +\frac{L\eta_{t}^{2}}{2}\|g_{t}\|^{2}.
    \tag{2.1}
\]

\noindent\textbf{Step 3 (Take conditional expectation).}  
Condition on $\mathcal{F}_{t}$ and use unbiasedness:
\[
    \E\!\bigl[f(x_{t+1})\mid\mathcal{F}_{t}\bigr]
    \le f(x_{t}) -\eta_{t}\|\nabla f(x_{t})\|^{2}
        +\frac{L\eta_{t}^{2}}{2}\E\!\bigl[\|g_{t}\|^{2}\mid\mathcal{F}_{t}\bigr].
    \tag{3.1}
\]

\noindent\textbf{Step 4 (Bound the second moment).}  
From the variance bound,
\[
    \E\!\bigl[\|g_{t}\|^{2}\mid\mathcal{F}_{t}\bigr]
    =\|\nabla f(x_{t})\|^{2}+ \E\!\bigl[\|g_{t}-\nabla f(x_{t})\|^{2}\mid\mathcal{F}_{t}\bigr]
    \le \|\nabla f(x_{t})\|^{2}+\sigma^{2}.
    \tag{4.1}
\]

\noindent\textbf{Step 5 (Insert adaptive step size).}  
Recall $\eta_{t}= \frac{\eta}{1+\gamma\|\widehat v_{t}\|}$. Since $\|\widehat v_{t}\|\ge0$ we have
\[
    \eta_{t}\le\eta,\qquad
    \eta_{t}^{2}\le\frac{\eta^{2}}{(1+\gamma\|\widehat v_{t}\|)^{2}}
    \le\frac{\eta^{2}}{1+\gamma\|\widehat v_{t}\|}.
    \tag{5.1}
\]

\noindent\textbf{Step 6 (Combine 3.1–5.1).}  
Using (4.1) and (5.1) in (3.1):
\[
\begin{aligned}
    \E\!\bigl[f(x_{t+1})\mid\mathcal{F}_{t}\bigr]
    &\le f(x_{t}) -\eta_{t}\|\nabla f(x_{t})\|^{2}
        +\frac{L\eta^{2}}{2(1+\gamma\|\widehat v_{t}\|)}\bigl(\|\nabla f(x_{t})\|^{2}+\sigma^{2}\bigr)\\
    &= f(x_{t}) -\underbrace{\Bigl(\eta_{t}-\frac{L\eta^{2}}{2(1+\gamma\|\widehat v_{t}\|)}\Bigr)}_{\displaystyle\Delta_{t}}\|\nabla f(x_{t})\|^{2}
        +\frac{L\eta^{2}\sigma^{2}}{2(1+\gamma\|\widehat v_{t}\|)} .
\end{aligned}
\tag{6.1}
\]

\noindent\textbf{Step 7 (Lower bound $\Delta_{t}$).}  
Because $\eta\le 1/L$ (Assumption \ref{ass:param}), we have $L\eta\le1$ and thus
\[
    \frac{L\eta^{2}}{2(1+\gamma\|\widehat v_{t}\|)}\le\frac{\eta}{2(1+\gamma\|\widehat v_{t}\|)}.
\]
Consequently,
\[
    \Delta_{t}\ge \frac{\eta}{2(1+\gamma\|\widehat v_{t}\|)} .
    \tag{7.1}
\]

\noindent\textbf{Step 8 (Apply PL condition).}  
From Assumption \ref{ass:pl},
\[
    \|\nabla f(x_{t})\|^{2}\ge 2\mu\bigl(f(x_{t})-f^{\star}\bigr).
    \tag{8.1}
\]

\noindent\textbf{Step 9 (One‑step recursion).}  
Insert (8.1) and (7.1) into (6.1):
\[
\begin{aligned}
    \E\!\bigl[f(x_{t+1})-f^{\star}\mid\mathcal{F}_{t}\bigr]
    &\le f(x_{t})-f^{\star}
        -\frac{\eta\mu}{1+\gamma\|\widehat v_{t}\|}\bigl(f(x_{t})-f^{\star}\bigr)
        +\frac{L\eta^{2}\sigma^{2}}{2(1+\gamma\|\widehat v_{t}\|)}\\
    &=\Bigl(1-\frac{\eta\mu}{1+\gamma\|\widehat v_{t}\|}\Bigr)\bigl(f(x_{t})-f^{\star}\bigr)
        +\frac{L\eta^{2}\sigma^{2}}{2(1+\gamma\|\widehat v_{t}\|)} .
\end{aligned}
\tag{9.1}
\]

\noindent\textbf{Step 10 (Take full expectation).}  
Define $A_{t}:=1+\gamma\|\widehat v_{t}\|$. Using $\E[\,\cdot\,]=\E[\E[\,\cdot\,|\mathcal{F}_{t}]]$ and Lemma~\ref{lem:boundG} to bound $\E[1/A_{t}]\ge 1/(1+\gamma G)$, we obtain
\[
\begin{aligned}
    \E\!\bigl[f(x_{t+1})-f^{\star}\bigr]
    &\le \E\!\Bigl[\Bigl(1-\frac{\eta\mu}{A_{t}}\Bigr)(f(x_{t})-f^{\star})\Bigr]
        +\frac{L\eta^{2}\sigma^{2}}{2}\E\!\Bigl[\frac{1}{A_{t}}\Bigr]\\
    &\le \Bigl(1-\frac{\eta\mu}{1+\gamma G}\Bigr)\E\!\bigl[f(x_{t})-f^{\star}\bigr]
        +\frac{L\eta^{2}\sigma^{2}}{2(1+\gamma G)} .
\end{aligned}
\tag{10.1}
\]

\noindent\textbf{Step 11 (Unroll the recursion).}  
Let $\rho:=\dfrac{\eta\mu}{1+\gamma G}\in(0,1)$. Recurrence (10.1) yields by induction:
\[
    \E\!\bigl[f(x_{T})-f^{\star}\bigr]
    \le (1-\rho)^{T}\bigl(f(x_{0})-f^{\star}\bigr)
        +\frac{L\eta^{2}\sigma^{2}}{2(1+\gamma G)}\sum_{t=0}^{T-1}(1-\rho)^{t}.
\tag{11.1}
\]

\noindent\textbf{Step 12 (Geometric sum).}  
Since $\sum_{t=0}^{T-1}(1-\rho)^{t}\le \frac{1}{\rho}$,
\[
    \frac{L\eta^{2}\sigma^{2}}{2(1+\gamma G)}\sum_{t=0}^{T-1}(1-\rho)^{t}
    \le \frac{L\eta^{2}\sigma^{2}}{2(1+\gamma G)}\cdot\frac{1}{\rho}
    =\frac{\eta\sigma^{2}}{2\mu(1+\gamma G)} .
\tag{12.1}
\]

\noindent\textbf{Step 13 (Conclusion).}  
Plugging (12.1) into (11.1) gives exactly the bound stated in Theorem~\ref{thm:linear}. $\square$

\section*{Rate Derivation (With Constants)}
From Lemma~\ref{lem:boundG} we can take a concrete bound on $G$:
\[
    G\le \frac{\alpha}{1-\alpha}\bigl(LD+\sigma\bigr),\qquad
    D:=\sup_{t}\E\!\bigl[\|x_{t}-x^{\star}\|\bigr].
\]
Hence
\[
    \rho = \frac{\mu\eta}{1+\gamma G}
    \ge \frac{\mu\eta}{1+\gamma\frac{\alpha}{1-\alpha}(LD+\sigma)} .
\]
Choosing $\eta=\frac{1}{L}$ maximises the numerator while respecting Assumption~\ref{ass:param}. The resulting linear rate is
\[
    (1-\rho)^{T}= \Bigl(1-\frac{\mu/L}{1+\gamma\frac{\alpha}{1-\alpha}(LD+\sigma)}\Bigr)^{T}.
\]

\section*{Special Cases}
\begin{enumerate}
    \item \textbf{Exact gradients ($\sigma=0$).} The steady‑state error term disappears and we obtain pure linear convergence.
    \item \textbf{No EMA adaptation ($\alpha\to0$).} Then $G\to0$, $\rho\to\mu\eta$ and the bound reduces to the classic GD rate $1-\mu\eta$.
    \item \textbf{Heavy‑tailed noise with large $\sigma$.} The term $\frac{\eta\sigma^{2}}{2\mu(1+\gamma G)}$ dominates; increasing $\gamma$ (more aggressive step‑size decay) reduces the asymptotic error.
\end{enumerate}

\end{document}